\begin{table}[H]
\centering
\scriptsize
\setlength{\tabcolsep}{4pt}
\renewcommand{\arraystretch}{1.15}
\caption{Mordell Lean 4 项目模块结构与统计数据}
\label{tab:mordell-module-stats}
\begin{tabular}{@{}>{\raggedright\arraybackslash}p{0.29\textwidth}>{\raggedright\arraybackslash}p{0.43\textwidth}rr@{}}
\toprule
\textbf{文件} & \textbf{目的} & \textbf{LOC} & \textbf{声明} \\
\midrule
\texttt{ZsqrtdBasic.lean} &
二次整数环 \(\mathbb Z[\sqrt d]\) 与有理二次代数的基础接口；整理范数、迹、基和有限模结构。 &
177 & 17 \\
\texttt{Approximation.lean} &
有理数取整逼近和范数估计；为 \(d=-5,-6,-13\) 的类群代表约化提供计算引理。 &
327 & 11 \\
\texttt{Dedekind.lean} &
证明 \(\mathbb Z[\sqrt d]\) 在相应同余与无平方因子条件下具有整环和 Dedekind 整环结构。 &
306 & 14 \\
\texttt{ClassGroupApprox.lean} &
形式化类群代表元约化准则，将逼近结果转化为关于理想类代表的有限性控制。 &
139 & 8 \\
\texttt{Descent.lean} &
形式化 Mordell 下降法主线：因式分解、理想互素、类群下降、单位处理和参数化结论。 &
549 & 26 \\
\texttt{EuclideanNegTwo.lean} &
构造 \(\mathbb Z[\sqrt{-2}]\) 的显式除法算法，并给出 Euclidean domain 实例。 &
176 & 18 \\
\texttt{ClassNumberSmall.lean} &
验证 \(d=-1,-2,-5,-6,-13\) 所需的小类数条件，并给出若干具体 cube-root 输入。 &
868 & 44 \\
\texttt{Extensions.lean} &
将下降法结论包装为更易复用的解集等价命题和无解判别准则。 &
104 & 6 \\
\texttt{Solutions.lean} &
给出具体 Mordell 方程的最终结论：\(d=-1,-2,-13\) 的整数解以及 \(d=-5,-6\) 的无解性。 &
188 & 12 \\
\texttt{Basic.lean} &
聚合导入基础证明模块，作为项目内部的统一入口。 &
5 & 0 \\
\texttt{Mordell.lean} &
导出项目主要模块，作为 Lake 默认目标。 &
5 & 0 \\
\midrule
\textbf{合计} & \textbf{11 个 Lean 文件} & \textbf{2844} & \textbf{156} \\
\bottomrule
\end{tabular}
\end{table}
